\section{Framework Design}

\subsection{Design Motivation}

Existing agent benchmarks suffer from three fundamental limitations: (1)~\emph{task isolation}---each task is independent, preventing measurement of cross-task dependency effects; (2)~\emph{outcome-only evaluation}---only final scores are reported, with no visibility into process cost, failure patterns, or recovery behavior; (3)~\emph{infrastructure coupling}---each benchmark requires custom evaluation infrastructure, making cross-benchmark comparison and model-backend matrix evaluation impractical.

\eva{} addresses these limitations through a unified evaluation architecture that decouples workload definition from execution infrastructure, supports heterogeneous model providers and agent frameworks through a standardized API layer, and records both outcome and process metrics as first-class analytical objects.

\subsection{Unified Evaluation Framework}

\eva{}'s evaluation framework consists of five major components (Figure~\ref{fig:architecture}):

\textbf{Task Orchestrator.}
Responsible for workload scheduling, dependency resolution, resurrection triggering, and evaluation pipeline management.
The orchestrator maintains the serial task chain state, monitors per-task pass/fail status, and triggers resurrection when a task falls below the 60\% passing threshold.

\textbf{Agent Runtime Layer.}
Supports multiple agent frameworks through a unified interface, each representing a distinct interaction paradigm:

OpenHands SDK: Full agent framework with unlimited self-loop iteration, skill injection, and MCP tool integration.
Claude Code: File I/O and command execution agent with multi-turn conversation.
Codex CLI: Autonomous coding agent with sandboxed execution.
Kimi CLI: Conversational coding agent with plan management and background task execution.

All backends share a common task context format and produce standardized result structures, enabling direct cross-backend comparison.

\textbf{Model Access Layer (\textsc{AIhub}).}
A unified API gateway (\url{https://aihub.arcsysu.cn}) providing standardized OpenAI-compatible access to 21 large language models spanning proprietary and open-weight families.
This design eliminates API-specific integration overhead and ensures that evaluation differences reflect model and scaffold capabilities rather than infrastructure variability.

\textbf{Execution Environment.}
Each evaluation run operates in an isolated workspace containing a complete copy of the target repository with all dependencies (ANTLR, LLVM 14, pybind11) pre-installed.
For \yatcc{}-Hard, each run additionally launches a fresh container environment that is destroyed after completion, ensuring zero cross-run contamination.
The workspace is configured via CMake with a fixed toolchain; scoring scripts are deterministic.

\textbf{Observation \& Leaderboard Layer.}
Collects and aggregates evaluation results, computes utility-oriented metrics, and maintains dynamic leaderboards.
All raw evaluation data (scores, logs, process metrics) is preserved for reproducibility.

% TODO: Add Figure showing the five-component architecture

\subsection{Workloads: Compiler Construction as Realistic Evaluation}

\eva{} grounds its evaluation in compiler construction, a domain chosen for its strict causal dependencies, well-defined intermediate artifacts, and objective automated scoring.
Compiler engineering provides an ideal evaluation substrate because: (1)~the pipeline has \emph{strict causal dependencies}---each stage consumes the output of its predecessor; (2)~intermediate artifacts are \emph{well-defined and verifiable}---token streams, parse trees, IR modules, and assembly files have clear structural specifications; (3)~the project includes \emph{automated scoring scripts} with comprehensive test suites.

We provide two complementary workloads spanning a difficulty spectrum:

\subsubsection{\yatcc{}: Scaffolded Implementation}

The original \yatcc{} workload~\cite{yatcc2024} organizes six tasks along the compiler construction pipeline (Table~\ref{tab:tasks}).
Each task provides \emph{scaffolded starter code}---the agent fills in missing implementation details within a provided framework.
This workload measures an agent's ability to \emph{complete} a partially-implemented system.

\begin{table*}[t]
\centering
\small
\begin{tabular}{clll}
\toprule
\textbf{Task} & \textbf{Name} & \textbf{Objective} & \textbf{Input Artifact} \\
\midrule
0 & Environment Setup & Build \& verify toolchain & --- \\
1 & Lexer & Tokenize C source code & Preprocessed source \\
2 & Parser & Construct AST/ASG & Token stream (from T1) \\
3 & IR Generation & Emit LLVM IR & AST/JSON (from T2) \\
4 & IR Optimization & Implement LLVM Pass & Raw IR (from T3) \\
5 & Code Generation & Produce RV64 assembly & Optimized IR (from T4) \\
\bottomrule
\end{tabular}
\caption{\eva{} task chain. Each task consumes the output artifact of its predecessor, forming a strict serial dependency.}
\label{tab:tasks}
\end{table}

\subsubsection{\yatcc{}-Hard: Full From-Scratch Implementation}

\yatcc{}-Hard removes all code logic from the original tasks, retaining only file dependency relationships to ensure automated scripts remain functional.
Agents receive \emph{no compiler construction knowledge} and \emph{no implementation framework}---they must implement the full compiler from scratch.
The task descriptions are reduced to experiment requirements and scoring script usage instructions.
This workload measures an agent's ability to \emph{construct} a complete system from minimal specification.

The difficulty gap between \yatcc{} and \yatcc{}-Hard quantifies the value of scaffolding in agentic coding tasks and reveals which models depend on provided structure versus which can operate autonomously.

Importantly, these workloads have been validated through years of practical use in compiler construction education at Sun Yat-sen University, ensuring their authenticity as real engineering tasks rather than synthetic benchmarks.

\subsection{Resurrection Protocol}

The \resurrection{} mechanism addresses the fundamental evaluation challenge of serial workloads: when an agent fails at Task~$k$, downstream tasks become unreachable, and the benchmark can only report ``failure at $k$'' without diagnosing whether the agent could succeed at Tasks~$k+1, \ldots, 5$ given correct prerequisites.

\textbf{Definition.}
When Task~$k$ scores below the passing threshold (60\%), \resurrection{} is triggered:
(1)~the agent's failed artifact is replaced with the \emph{golden intermediate artifact}---the correct output produced by the reference implementation;
(2)~the build system is reconfigured via \texttt{config.cmake} (\texttt{TASK\_k\_REVIVE=ON});
(3)~the agent continues to Task~$k+1$ with the corrected state.

\textbf{Design rationale.}
Resurrection is an \emph{evaluation protocol}, not a scoring adjustment.
It decomposes the compound failure ``cannot reach Task~$k+1$'' into two independent measurements:
(1)~\emph{predecessor failure}: the agent could not produce a passing artifact at Task~$k$;
(2)~\emph{downstream capability}: given correct prerequisites, can the agent succeed at Task~$k+1$?

This decomposition enables node-level pass rate reporting even when the full pipeline is broken, significantly improving diagnostic resolution.
Agents are \emph{not informed} about resurrection---the mechanism is triggered externally by the evaluation orchestrator, eliminating moral hazard.

\subsection{Intermediate Artifacts and Verification}

Each task produces a well-defined intermediate artifact that serves as both the evaluation object for the current task and the input dependency for the next:
Task~1 produces a \emph{token stream}; Task~2 produces an \emph{AST/ASG}; Task~3 produces \emph{LLVM IR}; Task~4 produces \emph{optimized IR}; Task~5 produces \emph{RV64 assembly}.
Scoring is performed by dedicated test suites that execute the agent's code against comprehensive test cases.
Each task's score ranges from 0 to 100, with a passing threshold of 60\%.